\documentclass[11pt, twocolumn]{article}
%Packages 
\usepackage{graphicx} % Required for inserting images
\usepackage[english]{babel}
\usepackage{lipsum} % This package generates filler text
\usepackage{newtxtext} % Sets the default text font to Times New Roman
\usepackage{setspace} % Used to set line spacing
\usepackage{hyperref} % For adding links
\usepackage{amsmath} % For advanced math formatting
%\usepackage{natbib} % For better bibliography management
\usepackage[margin=0.5in]{geometry} % Increase margins to make the text area narrower
\usepackage[framemethod=TikZ]{mdframed} 
\usepackage[backend=biber]{biblatex}
\addbibresource{references.bib}


\onehalfspacing % Sets line spacing to double

% Manually defining the title 
\newcommand{\paperTitle}{XXXXXTitleXXXXXXX \\ XXXXX}


\newcommand{\authors}{
\begin{flushleft}

% Manually defining the authors/collaboraters
\begin{center} 
\Large{FIRST LAST NAME, XXXXX, XXXXXX, Dr.Mario Diaz}
\end{center}

\textit{The University of Texas Rio Grande Valley, Brownsville, TX 78520, USA} \\
\textit{The Department of Physics \& Astronomy} \\
\textit{Research Experience for Undergraduates} 
\textit{Summer 2023} \\
\textit{Email: \href{mailto:your.email@example.com}{your.email@example.com}}
%type your email}
\end{flushleft}
}

% Create a new environment that centers the title of the abstract

\newenvironment{centeredabstract}
{\begin{mdframed}[hidealllines=true,backgroundcolor=gray!20,linewidth=2pt,
  topline=true,
  bottomline=true]\begin{quote}\textbf{Abstract}}
{\end{quote}\end{mdframed}}

\begin{document}

\twocolumn[
% Display the title and authors
\begin{center}
{\LARGE \textbf{\paperTitle}} % Large, bold title
\end{center}
\vspace{.05cm} % Some space between the title and author details
{\large \authors} % Large author details

\vspace{0.05cm} % Some space between the authors and abstract

\begin{centeredabstract}
    Leave it at the end: A brief summary of the research paper that provides an overview of the study's objectives, methods, results, and conclusions... 
    Déjelo al final.: Un breve resumen del artículo de investigación que brinda una visión general de los objetivos, métodos, resultados y conclusiones del estudio....

    \noindent \textbf{Keywords:} Sun: abundances; stars: early-type; Galaxy: structure; galaxies: individual (M31)
\end{centeredabstract}
]

\section{Introduction}

Provide a brief introduction to the research topic, explaining its significance and the research question or objective...

Introducción: La sección de introducción proporciona información de fondo y contexto sobre el tema de investigación, establece el problema o pregunta de investigación y describe el propósito del estudio... 

notes: if you want to add another section under introduction  get rid of % symbol 
%\subsection{title}   

notas: si deseas agregar otra sección debajo de Introduccion, elimina el símbolo de porcentaje (%).
%\subsection{Título} 

To correctly reference your information, ensure you add the relevant BibTeX entry into your references.bib file.

Consider the following example: "According to Albert Einstein, the Moon is round\cite{sample1}." The citation key 'sample1' corresponds to a specific entry in the references.bib file. To find the details of 'sample1', refer to the references.bib file. 

Para referenciar correctamente tu información, asegúrate de agregar la entrada BibTeX correspondiente en tu archivo references.bib.

Considera el siguiente ejemplo: "Según el Albert Einstein, la Luna es redonda\cite{sample1}." La clave de citación 'sample1' corresponde a una entrada específica en el archivo references.bib. Para encontrar los detalles de 'sample1', consulta el archivo references.bib.


\section{Literature Review}
(OPTIONAL) A concise review of relevant literature and previous research that provides context for the research project and highlights any gaps or limitations in the existing knowledge.... 

(OPCIONAL) Revisión de literatura: Una revisión de la literatura relevante e investigaciones previas sobre el tema, resaltando los conocimientos existentes, teorías y brechas que el estudio actual busca abordar.
Metodología: Descripción detallada del diseño de investigación, métodos y procedimientos utilizados para recolectar y analizar datos. Esta sección debe incluir información sobre la muestra, técnicas de recolección de datos, herramientas o instrumentos utilizados y cualquier método estadístico o analítico aplicado.

notes: if you want to add another section under literature get rid of % symbol 
%\subsection{title}   

notas: si deseas agregar otra sección debajo de literatura, elimina el símbolo de porcentaje (%).
%\subsection{Título}  %Optional (opcional)
\section{Methods}
Describe the methodology or approach used in the research project. This section should provide enough detail for readers to understand how the study was conducted... 

Metodología: Descripción detallada del diseño de investigación, métodos y procedimientos utilizados para recolectar y analizar datos. Esta sección debe incluir información sobre la muestra, técnicas de recolección de datos, herramientas o instrumentos utilizados y cualquier método estadístico o analítico aplicado... 

notes: if you want to add another section under Methods get rid of % symbol 
%\subsection{title}   

notas: si deseas agregar otra sección debajo de metodologia, elimina el símbolo de porcentaje (%).
%\subsection{Título} 


\section{Results}

Presentation and interpretation of the findings obtained from the data analysis. This section may include figures, tables, or graphs to support the results.

Resultados: Presentación e interpretación de los hallazgos obtenidos del análisis de datos. Esta sección puede incluir tablas, gráficos o diagramas para respaldar los resultados... 

notes: if you want to add another section under results get rid of % symbol 
%\subsection{title}  

notas: si deseas agregar otra sección debajo de Resultados, elimina el símbolo de porcentaje %.
%\subsection{Título} 


%Example table 
\begin{table}[!htbp]
	\centering
	\begin{tabular}{c c c}
		\hline
		\hline
		Time [JD - 2457984] & Magnitude & $\sigma$ [mag] \\
		\hline
		0.675868 & 18.9 & 0.1 \\
		0.681551 & 19.12 & 0.09 \\
		\hline
		\hline
	\end{tabular}
	\caption{Example how to make a table Time series photometry}
	\label{tab:time_series_photometry_g}
\end{table}

%Example to insert images

\begin{figure}
 \centering
 \includegraphics[width=0.40\textwidth]{Images/ctmo-logo.png}
 \caption{The font size in the text and numerical values on the axes is similar to the font size of this heading.}
    \label{fig:ctmo}
\end{figure}
\section{Discussion}

Analyze and interpret the results in relation to the research question or objective. Discuss any patterns, trends, or unexpected findings and their implications.

Discusión: Interpretación y análisis de los resultados en relación con la pregunta de investigación o hipótesis. También se deben discutir las implicancias, limitaciones y áreas potenciales para investigaciones futuras. 

notes: if you want to add another section under Discussion get rid of % symbol 
%\subsection{title}   

notas: si deseas agregar otra sección debajo de Discusion, elimina el símbolo de porcentaje (%).
%\subsection{Título}  


\section{Conclusion}
Summarize the main findings of the research project and their significance. Briefly restate the research question or objective and discuss any potential avenues for future research. 

Conclusiones: Un resumen de los principales hallazgos del estudio, enfatizando su importancia y relevancia. Las conclusiones deben abordar directamente la pregunta de investigación y resaltar las contribuciones realizadas al campo... 





\section{Acknowledgments}
 (if necessary): A section to acknowledge any individuals or institutions that contributed to the research project or provided support during the REU or any other program.

(si es necesario): Una sección para reconocer a las personas o instituciones que contribuyeron al proyecto de investigación o brindaron apoyo durante el programa REU o cualquier otro programa. 

 %if necessary (si es necesario) 

\appendix 
\appendix 
\section{Appendix} %change title 

Appendices (if necessary): Supplementary materials such as additional data, questionnaires, interview transcripts, or any other relevant information that supports the main content of the paper...

Apéndices (si es necesario): Materiales complementarios como datos adicionales, cuestionarios, transcripciones de entrevistas u otra información relevante que respalde el contenido principal del artículo... %if necessary (si es necesario) 

\printbibliography

\end{document}

